% ===== §6 =====
\section{The Easing of a Class of Conflict}
\label{sec:conflict}

\noindent
A third work of relational understanding is the easing of conflict, and here the established knowledge is empirical and strong. Decades of observational research on couples, associated above all with Gottman~\citep{gottman1999marriage}, have identified patterns that distinguish bonds that endure from those that dissolve. The research marks certain corrosive modes of exchange, contempt and defensiveness and the rest, and certain reparative ones, and it predicts, from the observed ratio of these, the trajectory of a bond with a reliability that the study of intimate life rarely attains. This is a body of findings the present account has no wish to contest, and it is important to be clear about the relation between the two, since the account offers something of a different kind.

The present account does not compete with this research at the level of evidence, and it advances no observational findings of its own. What it offers is a mechanism, an account of why a large class of intimate conflict arises and why understanding of the kind described here eases it. The offer is complementary. The empirical research establishes what patterns of conflict and repair predict the fate of a bond; the present account proposes, for one important class of conflict, an explanation of what that conflict is at its root.

The proposal follows from the bridge already drawn. If each partner is the bearer of a symbolic world, and if the closeness of intimacy generates the illusion that the two inhabit a single world, then a characteristic form of conflict arises when the partners use the same words to render aspects that do not coincide, and each takes the other's divergent rendering as a fault, an unreasonableness, a bad will, and not as a difference of worlds. A dispute that presents itself as a disagreement about a matter, about whether one of them acted wrongly, about what was owed, about what a word like respect or space or family requires, is in a large class of cases a collision of two renderings that the partners have not recognized as two, each arguing from within a world the other has not entered and each hearing the other's words in a world the other does not inhabit. The conflict is real, and its felt content, the sense of being wronged or unheard, is real, and yet its root is not the matter it appears to concern. Its root is an unrecognized difference of symbolic worlds.

Here the relational character does its work. To ease a conflict of this class is not, in the first place, to resolve the matter it appears to concern, to adjudicate who was right about the thing disputed. It is to move the conflict from the level of content to the level of worlds, to recognize that the partners have been rendering a shared matter in divergent ways, and to enter one another's rendering far enough that each sees how the other's words, which sounded like unreason from within one's own world, are the natural speech of another. What eases the conflict is the crossing that the closeness of intimacy had made to seem unnecessary. Once each has entered the other's rendering, the matter that seemed a contest of right and wrong is observed as a relation between two worlds, and the heat that belonged to the sense of being wronged by an unreasonable other subsides into the different and workable difficulty of two people who render a shared life in different terms.

The generative character marks the further point that this easing is not a return to a prior peace. When two partners cross a conflict of this class, they do not merely remove a disturbance and restore the state before it. They generate, in the crossing, an observation of the difference between their worlds that neither held before, and this observation becomes part of the shared reality of the bond, available to the next collision and lessening the ground from which it springs. A bond in which such crossings are repeatedly made is one in which the partners accumulate a shared understanding of the very structure of their differences, so that the easing of conflict is generative of a resource that the easing of the next conflict draws upon. This is the sense in which understanding eases conflict by generating, across the disputes, a growing observation, and not by settling them, of the two worlds whose unrecognized difference the disputes expressed.

The relation to the empirical research can now be stated. The observed patterns of corrosion and repair that predict the fate of a bond are, on this account, the surface at which the recognition or the non-recognition of the other's world becomes visible. Contempt is the affective form of a refusal to grant that the other's rendering has a world behind it; repair is the affective form of a return to the crossing. The research describes what these patterns do to the trajectory of a bond; the present account proposes what they are, and why the practice of relational understanding, which is the practice of crossing, bears on them at their root.
